\hypertarget{worked-examples}{%
\chapter{21. Worked examples}\label{worked-examples}}

A single confined thing, followed from the floor to a running shell.
Each tier is a real shipped file; the prose is the file's own. Read top
to bottom and you watch confinement accrete: the floor denies, the
template widens it four different ways, the leaf pins it.

\hypertarget{base-confined-the-floor-every-kennel-stands-on}{%
\section{\texorpdfstring{\texttt{base-confined} --- the floor every
kennel stands
on}{base-confined --- the floor every kennel stands on}}\label{base-confined-the-floor-every-kennel-stands-on}}

Template: base-confined --- the floor every kennel stands on. Inherits:
nothing --- this is the root of every confined template. Threat
catalogue: THREATS.md v0.4

The minimal confined posture: no\_new\_privs, setuid/setgid/setcap
denial, the constructed-\$HOME view, proxy-only egress, the
cloud-metadata + link-local invariant denies (RFC1918 stays reachable
--- see {[}net{]}), a curated environment, and a defence-in-depth
seccomp filter.

Deny-by-default across the board: exec.allow is EMPTY, which now denies
ALL execution (execution is deny-by-default like fs and net) --- so
base-confined cannot run anything and is not usable on its own. A
derived template/leaf adds the binaries it needs to \texttt{exec.allow};
the compiler then resolves and grants exactly the shared libraries those
binaries link (see the \texttt{{[}lib{]}} section below). Only enforced
sections appear here; sections the runtime does not implement
(dbus/x11/ptrace/signal) are deliberately omitted --- declaring them
would warn.

See docs/design/05-templates.md and templates/README.md.

\begin{kennelcode}
template_name = "base-confined"
\end{kennelcode}

\hypertarget{capabilities-framework-invariants}{%
\subsection{Capabilities (framework
invariants)}\label{capabilities-framework-invariants}}

\begin{kennelcode}
[cap]
no_new_privs = true        # PR_SET_NO_NEW_PRIVS, always; cannot be set false
bounding_set = []          # drop the entire capability bounding set
\end{kennelcode}

\hypertarget{execution-deny-by-default}{%
\subsection{Execution
(deny-by-default)}\label{execution-deny-by-default}}

exec.allow is empty, which DENIES ALL EXECUTION: a merely-readable file
is not executable. A derived template/leaf adds the binaries it needs;
\texttt{permissive-exec} (a \texttt{**} entry) is the explicit opt-out
that restores the open posture.

There is intentionally NO exec.deny list: under deny-by-default it would
be moot (sudo/su/pkexec/\ldots{} are not in any allowlist, so they never
run; and the deny\_* flags below + no\_new\_privs neuter setuid
escalation regardless). A deny list that enforces nothing is theatre, so
it is omitted rather than shipped as reassurance.

\begin{kennelcode}
[exec]
allow = []
\end{kennelcode}

\begin{kennelcode}
deny_setuid = true         # refuse setuid binaries at execve (framework invariant)
deny_setgid = true         # refuse setgid binaries (framework invariant)
deny_setcap = true         # refuse file-capability binaries
deny_writable = true       # refuse execution of files in writable paths
\end{kennelcode}

\begin{kennelcode}
path = ["/usr/bin", "/usr/local/bin", "/bin"]
\end{kennelcode}

\hypertarget{filesystem}{%
\subsection{Filesystem}\label{filesystem}}

The constructed view (fs.home, below): \$HOME is a fresh tmpfs into
which only granted paths are bound. Non-granted paths do not exist in
the view (T1.1).

\begin{kennelcode}
[fs]
# Read baseline: the curated base every dynamically-linked workload needs
# (§4.2 construction-by-absence). Only these subtrees are bound into the
# view; everything else under /usr (headers, source, /usr/local) is simply not
# present. The filesystem floor prevents data leakage; the real enforcement is
# at exec.allow|deny (Landlock FS_EXECUTE). Absent subtrees on this host (e.g.
# /usr/lib64 on usrmerge distros, /usr/libexec on Debian) are silently skipped.
read = [
    # The curated /usr base:
    "/usr/bin/**", "/usr/sbin/**",        # binary search path (Landlock gates exec)
    "/usr/lib/**",                        # shared libraries, locale, multiarch
    "/usr/lib64/**",                      # 64-bit lib path (Fedora/RHEL; skipped if absent)
    "/usr/libexec/**",                    # helper binaries (git-core etc; skipped if absent)
    "/usr/share/**",                      # arch-independent data (terminfo, zoneinfo, CAs, locale)
    # usrmerge compat:
    "/lib/**", "/lib64/**",
    # TLS root CAs + linker config (not under /usr):
    "/etc/ssl/**", "/etc/pki/**",
    "/etc/ld.so.conf", "/etc/ld.so.conf.d/**", "/etc/ld.so.cache",
    # update-alternatives symlink farm: many tools (awk→gawk, vi, editor, pager,
    # java, …) are reached via /usr/bin/<tool> -> /etc/alternatives/<tool> -> real
    # binary. Without /etc/alternatives in the view the symlink dangles and the tool
    # is "command not found". It is non-sensitive (symlinks only), bound read-only.
    "/etc/alternatives/**",
    # The libc/NSS files are SYNTHESISED per-kennel (scrubbed of host specifics),
    # not bound from the host: /etc/{hosts,resolv.conf,nsswitch.conf,services,
    # protocols,passwd,group,host.conf}. See kenneld::etc and §7.2.5.
    "/proc/self/**", "/proc/cpuinfo", "/proc/meminfo", "/proc/version",
    "/sys/devices/system/cpu/**",        # processor-count detection
]
\end{kennelcode}

\begin{kennelcode}
# Write baseline: none. The workload's writable space is the constructed $HOME and
# the private /tmp (below); a usable template or leaf grants its project write paths.
write = []
\end{kennelcode}

\begin{kennelcode}
# There is intentionally NO fs.deny list. The constructed view is deny-by-default:
# only fs.read/fs.write paths exist in it at all, so the host's ~/.ssh, ~/.aws,
# /etc/shadow, /dev/mem, … are simply absent — not present-but-denied. And a deny
# that fell *inside* a granted directory could not be enforced anyway (Landlock is
# allow-only and cannot subtract a single path). A long credential denylist here
# would be reassurance theatre; the real control is granting narrowly, not denying
# broadly. (T1.1/T2.1 are defended by the view's absence-by-default, not a denylist.)
\end{kennelcode}

\begin{kennelcode}
# The constructed $HOME (T1.1). HOME is /home/<user> (the masked [identity].user,
# default `kennel`); granted ~/ paths are bound beneath it; the Landlock ruleset
# backstops on the resolved view (§7.2.3/.5). The home root is writable by default,
# but it is a fresh tmpfs — anything written directly there is ephemeral. Persistence
# is opt-in per path via `persist` (or a writable ~/ grant), which binds the real host
# inode read-write beneath the home. Set `readonly = true` to suppress the default
# home-write grant (only write-granted ~/ paths stay writable then).
[fs.home]
shadow = true
\end{kennelcode}

\begin{kennelcode}
# Private /tmp (T1.1): a fresh tmpfs; the host /tmp is invisible.
[fs.tmp]
writable = true
size = "512M"
\end{kennelcode}

\begin{kennelcode}
# Procfs: PID namespace + hidepid=2 — the workload sees only its own tree (T1.1, T1.6).
[fs.proc]
hidepid = true
\end{kennelcode}

\begin{kennelcode}
# Minimal /dev: a constructed tmpfs with only these nodes bound in (T1.1, T1.6). Each
# granted node is also Landlock read/write/ioctl-able (§8.1).
[fs.dev]
allow = [
    "/dev/null", "/dev/zero", "/dev/random", "/dev/urandom",
    "/dev/tty", "/dev/pts/**",
]
\end{kennelcode}

\hypertarget{libraries}{%
\subsection{Libraries}\label{libraries}}

There is no library allowlist, by design. Execution is gated by Landlock
FS\_EXECUTE, which the kernel enforces only at execve(2): on the
allowlisted binary AND its dynamic loader (PT\_INTERP/ld.so, resolved by
the compiler from each exec.allow binary). Shared libraries are mmap'd
by the loader, which Landlock does NOT gate, so they load with READ
alone (07-3-exec) --- the fs.read grants above (/usr, /lib, /lib64)
already make them readable. A library therefore cannot be
``execute-gated'' under Landlock; the kennel makes no such claim rather
than ship an unenforceable filter. (To stop a workload running NEW
programs, that is exec.allow --- default-deny; to stop it loading code
in-process, the workload must not be an interpreter you handed arbitrary
input.)

\hypertarget{network}{%
\subsection{Network}\label{network}}

All egress funnels through a per-kennel SOCKS5/HTTP proxy; the cgroup
BPF denies direct connect() to anything but the proxy (fail-closed). The
proxy resolves names via the OS resolver and vets the answers against
the allowlist + the invariant denies (rebinding defence) --- there is no
configurable {[}net.dns{]}.

\begin{kennelcode}
[net]
mode = "constrained"
# Both families' proxy listeners are on by default in the proxied modes; a family
# is enabled iff its address resolves (no separate on/off flag).
# Listener address is computed from the kennel's <tag>/<ctx> (§7.3.2); override the
# host offset/port within the kennel's own subnet with proxy_listen_v4_address =
# "offset:port" (offset 1..=14, default "1:1080"). Project Kennel injects
# HTTPS_PROXY / HTTP_PROXY / ALL_PROXY pointing at the computed proxy address.
\end{kennelcode}

\begin{kennelcode}
# Invariant denies — the destinations that are NEVER a legitimate egress target,
# enforced deny-first by the proxy (even in `open` mode) and non-removable by any
# derived template or leaf (T1.6). Deliberately NARROW: cloud metadata (the SSRF
# crown jewel) and link-local only.
\end{kennelcode}

\begin{kennelcode}
# RFC1918 / CGNAT are intentionally NOT invariant denies. Making private space
# permanently unreachable is self-defeating — a kennel routinely needs a local dev
# server, a LAN database, an internal registry, or a corp service. In `constrained`
# mode (the default) nothing private is reachable anyway unless a `[[net.proxy.allow]]`
# names it; in `open` mode the operator has opted into arbitrary egress. Either way
# that is the policy author's call, not an immovable floor.
[[net.proxy.deny.invariant]]
cidr = "169.254.169.254/32"
reason = "cloud metadata IPv4 — never permitted from a kennel (mandatory invariant)"
[[net.proxy.deny.invariant]]
cidr = "fd00:ec2::254/128"
reason = "AWS IPv6 metadata — never permitted"
[[net.proxy.deny.invariant]]
cidr = "fe80::/10"
reason = "IPv6 link-local — no legitimate egress destination"
\end{kennelcode}

\begin{kennelcode}
# Bind: a wildcard bind (0.0.0.0 / ::) is rewritten to the kennel's private
# loopback address (so dev servers work but are reachable only from inside the
# kennel's address space); host loopback is not bindable; no privileged ports.
[net.bind]
inaddr_any_policy = "rewrite"
in6addr_any_policy = "rewrite"
allow_host_loopback_v4 = false
allow_host_loopback_v6 = false
min_port = 1024
\end{kennelcode}

\begin{kennelcode}
# Force IPV6_V6ONLY=1 so a dual-stack socket cannot escape the v4 rewrite (T1.6).
[net.ipv6]
force_v6only = true
\end{kennelcode}

\begin{kennelcode}
# Per-kennel egress audit log (one JSONL record per request; denies at higher
# level). Written by the proxy; persists across runs.
[net.audit]
log_path = "~/.local/state/kennel/<kennel>/network.jsonl"
level = "summary"
\end{kennelcode}

\hypertarget{af_unix-sockets}{%
\subsection{AF\_UNIX sockets}\label{af_unix-sockets}}

Default-deny; the constructed view contains only granted sockets.
Abstract- namespace sockets are denied categorically --- enforced
natively by Landlock scoping (ABI 6, §8.1), not the legacy
seccomp/AppArmor fallback. (T1.6.)

\begin{kennelcode}
[unix]
abstract = "deny"
\end{kennelcode}

\hypertarget{process-introspection}{%
\subsection{Process introspection}\label{process-introspection}}

The PID namespace hides every process outside the kennel; hidepid
backstops it. (Cross-boundary ptrace and signalling are already
prevented by the PID/user namespace --- there is no separate control the
runtime enforces, so none is declared. The advisory
{[}unsafe.ptrace{]}/{[}unsafe.signal{]} sub-sections exist to
\emph{express} intent but warn that scoping comes from PID-ns/seccomp;
this template needs neither. T1.6.)

\hypertarget{environment}{%
\subsection{Environment}\label{environment}}

The environment is SYNTHESISED, not filtered. The spawn clears the
inherited env entirely (env\_clear) and rebuilds it: the framework sets
HOME, PATH, USER, LOGNAME, SHELL (and forwards the caller's TERM);
\texttt{set} below adds fixed extras. There is no
\texttt{pass}/\texttt{deny} here --- a passthrough/denylist is fiction
when nothing is inherited in the first place (the host's secrets,
AWS\_\emph{, }\_TOKEN, \ldots{} never reach the kennel because the whole
env is wiped, not because they were named in a denylist).

\begin{kennelcode}
[env]
set = { TMPDIR = "/tmp", NO_PROXY = "" }
\end{kennelcode}

\hypertarget{seccomp-defence-in-depth}{%
\subsection{Seccomp (defence in depth)}\label{seccomp-defence-in-depth}}

Most enforcement is at higher layers (Landlock for files, cgroup BPF for
net); this filter denies a small set of syscalls with no legitimate use
here and a history of exploit-chain involvement. Default action: errno
(EPERM).

\begin{kennelcode}
[seccomp]
profile = "default"
deny = [
    "userfaultfd", "perf_event_open", "bpf",
    "process_vm_readv", "process_vm_writev",
    "kexec_load", "kexec_file_load",
    "mount", "umount", "umount2", "pivot_root",
    "swapon", "swapoff", "reboot",
    "init_module", "finit_module", "delete_module",
    "personality",
]
\end{kennelcode}

\hypertarget{resource-limits-off-by-default}{%
\subsection{Resource limits (off by
default)}\label{resource-limits-off-by-default}}

setrlimit(2) caps, applied in the seal. Nothing is set here; a derived
template or leaf opts in. Names are the short rlimit resources (nofile,
nproc, as, cpu, data, fsize, core, stack, memlock, \ldots); values are
``soft'' or ``soft:hard'', each a number (optional K/M/G) or
``unlimited''. Example:

\begin{kennelcode}
[ulimits]
\end{kennelcode}

\begin{kennelcode}
nofile = "8192"
\end{kennelcode}

\begin{kennelcode}
nproc  = "512"
\end{kennelcode}

\begin{kennelcode}
core   = "0"
\end{kennelcode}

\begin{kennelcode}
[signature]
algorithm = "sshsig"
key_id = "kennel-maint-2026"
signature = "...envelope elided in print..."
\end{kennelcode}

\hypertarget{interactive-template-the-floor-made-into-a-usable-human-shell}{%
\section{\texorpdfstring{\texttt{interactive\ (template)} --- the floor
made into a usable human
shell}{interactive (template) --- the floor made into a usable human shell}}\label{interactive-template-the-floor-made-into-a-usable-human-shell}}

Template: interactive --- a credible human-driven shell.

A confined but genuinely usable workstation shell: bash + the standard
coreutils, text tools, editors, archives, process inspection, and the
common network tools, with HOST egress (a human is at the keyboard;
direct egress on the host stack, the cloud-metadata invariant still
applies). Execution is deny-by-default --- the toolset below is exactly
what runs; the compiler resolves each tool's library closure
(\texttt{{[}lib{]}}, inherited).

\$HOME and /tmp are writable (ephemeral); grant a project dir +
persistence in your own leaf. For an offline shell, derive base-confined
instead (or set net.mode).

\begin{kennelcode}
# run it directly:
kennel run interactive -- /bin/bash
# or derive it and add your project:
#   template_base = "interactive"
#   [[fs.write.add]] path = "~/work"
\end{kennelcode}

\begin{kennelcode}
template_name = "interactive"
template_base = "base-confined"
\end{kennelcode}

Exec floor composed from the fragment catalogue (05-templates.md §5.10):
the shells + POSIX userland + capability bundles, instead of
hand-listing them.

\begin{kennelcode}
include = ["core-shell", "core-coreutils", "core-file-mutation", "core-archive", "net-clients", "vcs-git"]
\end{kennelcode}

Host egress: a human-driven shell needs unmediated network (curl, git
clone, ssh, \ldots) on the host stack --- no SOCKS proxy. This shares
the host network namespace and so reinstates the host-recon residual
(T1.6), which is why \texttt{host} requires a reason; the cloud-metadata
invariant deny + the net.bpf allowlist are still enforced via
BPF/Landlock.

\begin{kennelcode}
[net]
mode = "host"
reason = "an interactive human shell needs direct, unmediated egress (git/ssh/curl) on the host network stack"
\end{kennelcode}

\begin{kennelcode}
# Re-list base-confined's read baseline (a template's fs.read replaces, not adds) and
# widen /proc to the whole namespaced procfs so `ps`/`top`/`free` see this kennel's
# own processes (the PID namespace already hides everything else).
[fs]
read = [
    "/usr/**", "/lib/**", "/lib64/**",
    "/etc/ssl/**", "/etc/pki/**",
    "/etc/ld.so.conf", "/etc/ld.so.conf.d/**", "/etc/ld.so.cache",
    "/etc/alternatives/**",   # awk→gawk, vi, editor, pager … symlink farm (else dangles)
    "/proc/**",
    "/sys/devices/system/cpu/**",
]
\end{kennelcode}

\begin{kennelcode}
[exec]
shell = "/bin/bash"
# Only the binaries no fragment provides stay inline here; the shell, the POSIX
# userland, and the toolchains come from the `include` above.
allow = [
    "/usr/bin/clear",
    "/usr/bin/df",
    "/usr/bin/du",
    "/usr/bin/free",
    "/usr/bin/hexdump",
    "/usr/bin/ip",
    "/usr/bin/kill",
    "/usr/bin/nano",
    "/usr/bin/netstat",
    "/usr/bin/od",
    "/usr/bin/pgrep",
    "/usr/bin/pidof",
    "/usr/bin/pkill",
    "/usr/bin/ps",
    "/usr/bin/python3",
    "/usr/bin/scp",
    "/usr/bin/ss",
    "/usr/bin/ssh",
    "/usr/bin/strings",
    "/usr/bin/sync",
    "/usr/bin/timeout",
    "/usr/bin/top",
    "/usr/bin/uptime",
    "/usr/bin/vi",
    "/usr/bin/vim",
    "/usr/bin/vim.basic",
    "/usr/bin/watch",
    "/usr/bin/xxd",
    "/usr/sbin/arp",
    "/usr/sbin/ifconfig",
    "/usr/sbin/route",
]
# net-tools (ifconfig/route/arp) live in /usr/sbin; put it on PATH so they resolve
# by name. Landlock EXECUTE is independent of PATH (the absolute allow entries above
# are what grant execution); this is purely so `ifconfig` is found.
path = ["/usr/bin", "/usr/local/bin", "/bin", "/usr/sbin", "/sbin"]
\end{kennelcode}

\begin{kennelcode}
[lifecycle]
ttl = "12h"
ttl_action = "warn"
\end{kennelcode}

\begin{kennelcode}
[signature]
algorithm = "sshsig"
key_id = "kennel-maint-2026"
signature = "...envelope elided in print..."
\end{kennelcode}

\hypertarget{include-core-shell}{%
\subsection{\texorpdfstring{include\ =\ "core-shell"}{include = "core-shell"}}\label{include-core-shell}}

Fragment: core-shell (POSIX shells)

A composable capability bundle (05-templates.md §5.10): The POSIX
shells. base-confined denies all execution and sets
\texttt{{[}exec{]}.shell} without granting it, so every interactive or
scripted kennel needs the shell explicitly; this is that grant, once,
instead of in every template.

\begin{kennelcode}
include = ["core-shell"]
\end{kennelcode}

Additive-only (§5.10): every entry is an
\texttt{{[}{[}exec.allow.add{]}{]}} delta. Composing this widens the
program menu only; \texttt{argv{[}0{]}} stays gated by the resolved
\texttt{{[}exec{]}.allow} under Landlock execve default-deny, and the
cage (net/fs/ttl/ ceilings) is untouched.

Adds to the exec floor:

\begin{itemize}
\tightlist
\item
  \texttt{/bin/sh} (the POSIX shell, and base-confined's default
  \texttt{{[}exec{]}.shell})
\item
  \texttt{/usr/bin/sh} (the POSIX shell on a merged-usr system)
\item
  \texttt{/bin/bash} (GNU bash)
\item
  \texttt{/usr/bin/bash} (GNU bash on a merged-usr system)
\item
  \texttt{/usr/bin/dash} (the Debian Almquist shell (\texttt{/bin/sh} on
  Debian/Ubuntu))
\item
  \texttt{/usr/bin/env} (the
  \texttt{\#!/usr/bin/env\ \textless{}interp\textgreater{}} shebang
  interpreter (the kernel's execve target for env-shebang scripts))
\end{itemize}

\hypertarget{include-core-coreutils}{%
\subsection{\texorpdfstring{include\ =\ "core-coreutils"}{include = "core-coreutils"}}\label{include-core-coreutils}}

Fragment: core-coreutils (Read / compute / text userland)

A composable capability bundle (05-templates.md §5.10): The non-mutating
POSIX userland: read, inspect, filter, and transform. The block every
interactive and agent template repeats. Carries NO filesystem-mutating
tool --- compose \texttt{core-file-mutation} for that --- so a read-only
kennel includes this alone and cannot write.

\begin{kennelcode}
include = ["core-coreutils"]
\end{kennelcode}

Additive-only (§5.10): every entry is an
\texttt{{[}{[}exec.allow.add{]}{]}} delta. Composing this widens the
program menu only; \texttt{argv{[}0{]}} stays gated by the resolved
\texttt{{[}exec{]}.allow} under Landlock execve default-deny, and the
cage (net/fs/ttl/ ceilings) is untouched.

Adds to the exec floor:

\begin{itemize}
\tightlist
\item
  \texttt{/usr/bin/cat} (concatenate and print files)
\item
  \texttt{/usr/bin/ls} (list directory contents)
\item
  \texttt{/usr/bin/head} (print the first lines)
\item
  \texttt{/usr/bin/tail} (print the last lines)
\item
  \texttt{/usr/bin/wc} (count lines/words/bytes)
\item
  \texttt{/usr/bin/sort} (sort lines)
\item
  \texttt{/usr/bin/uniq} (filter adjacent duplicate lines)
\item
  \texttt{/usr/bin/cut} (extract columns)
\item
  \texttt{/usr/bin/tr} (translate characters)
\item
  \texttt{/usr/bin/nl} (number lines)
\item
  \texttt{/usr/bin/tac} (concatenate in reverse)
\item
  \texttt{/usr/bin/rev} (reverse lines)
\item
  \ldots and 36 more
\end{itemize}

\hypertarget{include-core-file-mutation}{%
\subsection{\texorpdfstring{include\ =\ "core-file-mutation"}{include = "core-file-mutation"}}\label{include-core-file-mutation}}

Fragment: core-file-mutation (Filesystem-mutating userland)

A composable capability bundle (05-templates.md §5.10): The write-side
coreutils: create, move, remove, link, change mode. Kept separate from
\texttt{core-coreutils} so a read-only kennel cannot mutate while a
scratch/build kennel adds it.

\begin{kennelcode}
include = ["core-file-mutation"]
\end{kennelcode}

Additive-only (§5.10): every entry is an
\texttt{{[}{[}exec.allow.add{]}{]}} delta. Composing this widens the
program menu only; \texttt{argv{[}0{]}} stays gated by the resolved
\texttt{{[}exec{]}.allow} under Landlock execve default-deny, and the
cage (net/fs/ttl/ ceilings) is untouched.

Adds to the exec floor:

\begin{itemize}
\tightlist
\item
  \texttt{/usr/bin/cp} (copy files)
\item
  \texttt{/usr/bin/mv} (move/rename files)
\item
  \texttt{/usr/bin/rm} (remove files)
\item
  \texttt{/usr/bin/mkdir} (create directories)
\item
  \texttt{/usr/bin/rmdir} (remove empty directories)
\item
  \texttt{/usr/bin/ln} (create links)
\item
  \texttt{/usr/bin/touch} (create/update timestamps)
\item
  \texttt{/usr/bin/chmod} (change file modes)
\item
  \texttt{/usr/bin/mktemp} (create a temporary file/dir)
\item
  \texttt{/usr/bin/install} (copy and set permissions)
\end{itemize}

\hypertarget{include-core-archive}{%
\subsection{\texorpdfstring{include\ =\ "core-archive"}{include = "core-archive"}}\label{include-core-archive}}

Fragment: core-archive (Archive + compression)

A composable capability bundle (05-templates.md §5.10): Tar and the
common compressors/archivers --- unpack a source tarball, compress an
artifact.

\begin{kennelcode}
include = ["core-archive"]
\end{kennelcode}

Additive-only (§5.10): every entry is an
\texttt{{[}{[}exec.allow.add{]}{]}} delta. Composing this widens the
program menu only; \texttt{argv{[}0{]}} stays gated by the resolved
\texttt{{[}exec{]}.allow} under Landlock execve default-deny, and the
cage (net/fs/ttl/ ceilings) is untouched.

Adds to the exec floor:

\begin{itemize}
\tightlist
\item
  \texttt{/usr/bin/tar} (tape archiver)
\item
  \texttt{/usr/bin/gzip} (gzip compress)
\item
  \texttt{/usr/bin/gunzip} (gzip decompress)
\item
  \texttt{/usr/bin/zcat} (stream a gzip file)
\item
  \texttt{/usr/bin/xz} (xz (de)compress)
\item
  \texttt{/usr/bin/bzip2} (bzip2 (de)compress)
\item
  \texttt{/usr/bin/zip} (zip archiver)
\item
  \texttt{/usr/bin/unzip} (zip extractor)
\item
  \texttt{/usr/bin/zstd} (zstandard (de)compress)
\end{itemize}

\hypertarget{include-net-clients}{%
\subsection{\texorpdfstring{include\ =\ "net-clients"}{include = "net-clients"}}\label{include-net-clients}}

Fragment: net-clients (HTTP(S) fetch clients)

A composable capability bundle (05-templates.md §5.10): The command-line
fetch-client BINARIES. Distinct from \texttt{net-permissive}, which
grants egress DESTINATIONS: this adds curl/wget to \texttt{exec.allow};
the leaf or \texttt{net-permissive} says where they may reach, and a
\texttt{{[}net{]}.mode\ =\ "constrained"} cage + the proxy bound that
egress.

\begin{kennelcode}
include = ["net-clients"]
\end{kennelcode}

Additive-only (§5.10): every entry is an
\texttt{{[}{[}exec.allow.add{]}{]}} delta. Composing this widens the
program menu only; \texttt{argv{[}0{]}} stays gated by the resolved
\texttt{{[}exec{]}.allow} under Landlock execve default-deny, and the
cage (net/fs/ttl/ ceilings) is untouched.

Adds to the exec floor:

\begin{itemize}
\tightlist
\item
  \texttt{/usr/bin/curl} (transfer a URL)
\item
  \texttt{/usr/bin/wget} (retrieve files over HTTP(S)/FTP)
\end{itemize}

\hypertarget{include-vcs-git}{%
\subsection{\texorpdfstring{include\ =\ "vcs-git"}{include = "vcs-git"}}\label{include-vcs-git}}

Fragment: vcs-git

A composable capability bundle (05-templates.md §5.10): git plus its
per-subcommand helpers, and the system git config read-only. Include it
for a kennel that does source control.

\begin{kennelcode}
include = ["vcs-git"]
\end{kennelcode}

Additive-only (§5.10). This grants the git \emph{binaries} and the
\emph{system} config; it deliberately does NOT grant network egress ---
git-over-HTTPS to a specific forge is a destination the leaf states (or
\texttt{net-permissive} adds), and git-over-SSH goes through the
per-kennel SSH bastion (\texttt{{[}ssh{]}}, §7.10), never a real key in
the kennel.

Adds to the exec floor:

\begin{itemize}
\tightlist
\item
  \texttt{/usr/bin/git} (git version-control client)
\item
  \texttt{/usr/lib/git-core/**} (git's per-subcommand helper binaries,
  re-exec'd by \texttt{git})
\item
  \texttt{/etc/gitconfig} (system-wide git configuration (read-only))
\end{itemize}

\hypertarget{interactive-policy-the-thin-leaf-that-pins-a-workload}{%
\section{\texorpdfstring{\texttt{interactive\ (policy)} --- the thin
leaf that pins a
workload}{interactive (policy) --- the thin leaf that pins a workload}}\label{interactive-policy-the-thin-leaf-that-pins-a-workload}}

Reference policy: interactive --- the permissive, ready-to-run confined
human shell.

Maintainer-signed and shipped enabled: \texttt{kennel\ run\ interactive}
drops you into a confined \texttt{/bin/bash} with a broad workstation
toolset --- the interactive base (shells, coreutils, file tools,
editors, archives, net clients, git, python3) widened with the C
toolchain and the development headers --- and HOST egress. Deliberately
permissive: a credible daily-driver shell, still inside the cage
(no\_new\_privs, dropped caps, deny-setuid, shadowed \$HOME, masked
/etc, private /tmp, self-only procfs). Derive your own leaf to add a
project dir or tighten the toolset.

\begin{kennelcode}
name = "interactive"
template_base = "interactive"
\end{kennelcode}

Widen the inherited interactive toolset toward a full dev shell: the C
toolchain (cc/make/\ldots) and the development headers, on top of the
base's shells/coreutils/editors/net-clients/git/python3. The
language-registry fragments (lang-python/lang-node) are NOT included ---
they carry by-name proxy-egress rules host mode rejects; under host
egress python3 (inherited) reaches registries directly.

\begin{kennelcode}
include = [
    "toolchain-c",
    "dev-headers",
]
\end{kennelcode}

Host egress (also the interactive default): a human shell needs
unmediated network.

\begin{kennelcode}
[net]
mode = "host"
reason = "an interactive human shell needs direct, unmediated egress (git/ssh/curl) on the host network stack"
\end{kennelcode}

\begin{kennelcode}
[workload]
argv = ["/bin/bash"]
\end{kennelcode}

\begin{kennelcode}
[signature]
algorithm = "sshsig"
key_id = "kennel-maint-2026"
signature = "...envelope elided in print..."
\end{kennelcode}

\hypertarget{include-toolchain-c}{%
\subsection{\texorpdfstring{include\ =\ "toolchain-c"}{include = "toolchain-c"}}\label{include-toolchain-c}}

Fragment: toolchain-c

A composable capability bundle (05-templates.md §5.10): the C/C++ build
toolchain --- compiler, assembler, linker, archiver, and make. Include
it for a kennel that compiles native code; the project tree itself is
granted by the leaf, so this fragment adds only the tools.

\begin{kennelcode}
include = ["toolchain-c"]
\end{kennelcode}

Additive-only (§5.10). No network and no extra filesystem write: a build
reads the toolchain (already on \texttt{fs.read} via base-confined's
\texttt{/usr/**}) and writes only into the project tree the leaf made
writable.

Adds to the exec floor:

\begin{itemize}
\tightlist
\item
  \texttt{/usr/bin/cc} (C compiler driver (the default \texttt{cc}))
\item
  \texttt{/usr/bin/gcc} (GNU C compiler)
\item
  \texttt{/usr/bin/g++} (GNU C++ compiler)
\item
  \texttt{/usr/bin/cpp} (C preprocessor (invoked by the compiler
  driver))
\item
  \texttt{/usr/bin/as} (GNU assembler)
\item
  \texttt{/usr/bin/ld} (GNU linker)
\item
  \texttt{/usr/bin/ar} (static-archive tool)
\item
  \texttt{/usr/bin/make} (GNU make)
\item
  \texttt{/usr/lib/gcc/**} (gcc's per-version backend binaries (cc1,
  cc1plus, collect2), re-exec'd by the driver)
\end{itemize}

\hypertarget{include-dev-headers}{%
\subsection{\texorpdfstring{include\ =\ "dev-headers"}{include = "dev-headers"}}\label{include-dev-headers}}

Fragment: dev-headers (Development headers + kernel source)

A composable capability bundle (05-templates.md §5.10): Adds /usr/src
and /usr/include back into the view (read-only). These are absent from
the curated base (construction-by-absence) and only needed by
build/compile workloads. Include this fragment in templates that compile
C/C++/Rust code against system headers.

\begin{kennelcode}
include = ["dev-headers"]
\end{kennelcode}

Additive-only (§5.10): every entry is an {[}{[}fs.read.add{]}{]} delta.
Composing this widens the read view only; the cage
(net/fs.write/exec/ceilings) is untouched.

Adds to the exec floor:

\begin{itemize}
\tightlist
\item
  \texttt{/usr/include/**} (C/C++ system headers (absent from the
  curated base))
\item
  \texttt{/usr/src/**} (kernel headers and source trees (absent from the
  curated base))
\end{itemize}
